% Complete finite proofs for P16-C1.
\section{Finite logical abduction: proofs}
\label{sec:annex-p16-finite-abduction}

Use the definitions of Chapter~\ref{sec:p16-finite-abduction}. The source
model and its single-deletion criterion are Eiter and Gottlob's
\cite[sec.~2, manuscript pp.~9--11]{EiterGottlobAbductionManuscript}.
The finite enumeration, extended inconsistent-input case and examples are
proved explicitly here.

\begin{annexitem}{annex:p16.abduction-semantics}{thm:p16-abduction-semantics}{Finite valuations and the explanation criterion}
\begin{proof}
\textbf{Agreement with the book's semantics.}\enspace
Let \(v:V\to\mathbb{B}\). Extend it to a valuation
\(\widetilde v:\mathsf{Prop}\to\mathbb{B}\) by giving every atom outside \(V\)
the value \(\mathbf{F}\). For each formula over \(V\), its recursively
computed value under \(v\) equals its value under \(\widetilde v\).
This follows by induction on the formula: the atomic values agree,
\(\bottom\) has the same fixed value, and each connective applies the
same truth function to the values of its immediate subformulas. Conversely,
restriction of any valuation on \(\mathsf{Prop}\) to \(V\) preserves the
values of all formulas over \(V\) by the same induction. This is also
the finite restriction instance of \cref{lem:prop-coincidence}.

It follows that \(\Gamma\) has a global satisfying valuation exactly
when \(\operatorname{Mod}_V(\Gamma)\) is nonempty. For a formula
\(\varphi\) over \(V\), \(\Gamma\entails\varphi\) holds
exactly when every member of \(\operatorname{Mod}_V(\Gamma)\)
satisfies \(\varphi\): extend a local countervaluation for one
direction and restrict a global countervaluation for the other. Apply
this to \(\Gamma=T\cup E\) and to each atom \(m\in M\).
An atom has value \(\mathbf{T}\) exactly when the valuation satisfies it,
which proves the finite criterion.

The source's theory is finite and consistent, its \(H,M\) are subsets
of the same finite variable alphabet, and its solution condition is exactly
consistency of \(T\cup E\) and consequence of every manifestation.
Its finite alphabet can be named by distinct members of \(\mathsf{Prop}\);
transporting valuations along that naming preserves atomic truth values.
The source's truth constant is translated to \(\neg\bottom\), which
is always true. Its biconditional is translated to the conjunction of the
two material implications. That conjunction is true exactly when the two
formula values agree: equal values make both implications true, and unequal
values make the implication from the true value to the false value false.
Induction on the source formula now proves preservation of truth under
these expansions, since the other constructors have the same truth clauses.
The naming sends subsets of hypotheses bijectively to subsets of the named
hypotheses and preserves both inclusion and cardinality. Thus the criterion
agrees on every source input under this stated notation translation.
The only additional input
case introduced in this chapter drops the consistency requirement on \(T\);
the next paragraph determines that case directly.

\paragraph{Boundary cases.}
Every model of \(T\cup E\) is a model of \(T\), so
\[
 \operatorname{Mod}_V(T\cup E)\subseteq\operatorname{Mod}_V(T).
\]
If the latter is empty, so is the former, for every candidate. Thus an
inconsistent background permits no explanation. In this case semantic
entailment is vacuous, but the distinct consistency conjunct is false.
Putting \(E=\varnothing\) gives \(T\cup E=T\), which proves the
empty-explanation equivalence. If \(M=\varnothing\), every valuation
satisfies every member of \(M\), because it has no members; explanation
then reduces precisely to consistency of \(T\cup E\). All arguments
also apply when \(V\) is empty: there is exactly one function from the
empty set to \(\mathbb{B}\).
\end{proof}
\end{annexitem}

\begin{annexitem}{annex:p16.abduction-minimality}{thm:p16-abduction-minimality}{Minimal explanations and single deletions}
\begin{proof}
\textbf{Qualified monotonicity.}\enspace
If \(E\subseteq F\), every valuation satisfying \(T\cup F\)
satisfies \(T\cup E\). Therefore, if \(E\) explains the observations,
every model of \(T\cup F\) satisfies \(M\). When \(T\cup F\)
is also consistent, this proves that \(F\) is an explanation. The
consistency condition cannot be inferred from \(E\subseteq F\).

\paragraph{Finite minimality.}
Fix \(E\in\operatorname{Expl}(\mathcal P)\). The family
\[
 \{F\subseteq E:F\in\operatorname{Expl}(\mathcal P)\}
\]
is finite and nonempty, since it contains \(E\). Select a member
\(E_0\) of least cardinality in this family. A proper subset of
\(E_0\) which were an explanation would also be a subset of \(E\)
and would have smaller cardinality. That contradicts the choice. Hence
\(E_0\) is subset-minimal among all explanations and is contained in
\(E\).

If the full explanation family is nonempty, it too is a finite subfamily
of \(2^H\), so it has a member of least cardinality. Every such member
is subset-minimal, because a proper explanatory subset would have smaller
cardinality. This proves existence of both families and the asserted
inclusions. If two distinct subset-minimal explanations were comparable,
one would be a proper explanatory subset of the other, contradicting the
latter's minimality. Their cardinalities need not agree: the complete
calculation for \(\mathcal P_{\rm sizes}\) below gives sizes one and two.
That calculation uses the definitions directly, without invoking this
incomparability assertion.

\paragraph{Single-deletion characterization.}
Let \(E\) be an explanation. If \(E\) is subset-minimal, every
\(E\setminus\{h\}\), for \(h\in E\), is a proper subset and
cannot be an explanation.

Conversely, suppose \(E\) is not subset-minimal. Choose an explanation
\(F\subsetneq E\) and some \(h\in E\setminus F\). Put
\(D=E\setminus\{h\}\). Then \(F\subseteq D\subsetneq E\).
A model of \(T\cup E\) exists and also satisfies \(T\cup D\),
so \(T\cup D\) is consistent. Every model of \(T\cup D\)
satisfies \(T\cup F\), and hence satisfies \(M\), because \(F\)
is an explanation. Thus \(D\) is an explanation, contradicting failure
of all one-atom deletions. This proves the converse, including the
consistency step in the source's criterion. If \(E=\varnothing\),
there is no proper subset and no atom to delete; both conditions hold.
If the background is inconsistent, there is no explanation \(E\) to
which the characterization would apply, so the extended convention adds
no exception to the statement.

Finally, the constraint example below gives an explanatory \(\{a\}\)
whose permitted superset \(\{a,b\}\) is inconsistent. It demonstrates
the stated limitation of unqualified upward closure by a direct calculation.
\end{proof}
\end{annexitem}

\begin{annexitem}{annex:p16.abduction-enumeration}{thm:p16-abduction-enumeration}{Soundness, completeness and finite termination of the construction}
\begin{proof}
Let the fixed lists of hypotheses and variables have lengths \(r\) and
\(n\), respectively. Each binary vector of length \(r\) picks a
unique subset of \(H\), and every subset has exactly its membership
vector. Each binary vector of length \(n\) likewise determines a unique
valuation on \(V\), and every valuation has exactly that vector. A
length-zero vector is unique, so these correspondences include the empty
cases. Induction on the length shows that the numbers of vectors are
\(2^r\) and \(2^n\): each extension by one position has the two
choices zero and one, starting from one empty vector.

For a fixed candidate \(E\), induction on the number of scanned
valuations gives two flag invariants. The first flag is true exactly when
one of the scanned valuations satisfies \(T\cup E\). The second is
true exactly when one of them satisfies \(T\cup E\) and falsifies
at least one \(m\in M\). Both invariants hold before any scan because
the flags are false and no valuation has been scanned. At the next
valuation, the two stated tests add exactly its corresponding witness, and
the flags retain all previous witnesses. This proves the induction step.

After the complete valuation scan, the first flag is true exactly when
\(\operatorname{Mod}_V(T\cup E)\) is nonempty, and the second is
false exactly when no model in this set falsifies an observation. By
\cref{thm:p16-abduction-semantics}, their tested combination holds exactly
when \(E\) is an explanation. Thus every appended set is an explanation
(soundness), and every explanation is appended when its unique candidate
vector is processed (completeness). No candidate is processed twice, and
each candidate is appended at most once. Hence
\(\mathcal L=\operatorname{Expl}(\mathcal P)\) with no duplicates.

The subset filter keeps an explanation \(E\) exactly when the complete
list contains no explanatory proper subset of it. Since the list is
complete, this is precisely the definition of
\(\operatorname{MinExpl}(\mathcal P)\); both directions follow from
that equivalence. If \(\mathcal L\) is empty, there is no cardinality
minimum to choose and the stated empty output agrees with its definition.
Otherwise a finite scan maintains the least cardinality seen so far,
starting with the cardinality of the first member. Comparing each
successive cardinality with this value preserves that invariant, so the
final value is the least over the entire list. Retaining exactly the
members of that size gives every minimum-cardinality explanation and only
those explanations. Neither filtering step duplicates members.

Every formula of \(T\) is a finite tree, so recursive truth evaluation
terminates by induction on that tree. The theory, candidate and observation
sets are finite, so the tests at a candidate--valuation pair terminate.
There are exactly \(2^r\cdot2^n=2^{r+n}\) such pairs in the prescribed
unshortened scans. Each of the two filtering procedures scans finite lists
and finite sets. Consequently the whole construction terminates, including
inconsistent theories and empty inputs. This proves a bound on the stated
pair inspections, without replacing formula evaluation or set comparison
by unit-cost assumptions.
\end{proof}
\end{annexitem}

\begin{annexitem}{annex:p16.abduction-examples}{thm:p16-abduction-examples}{Complete calculations for all six inputs}
\begin{proof}
\textbf{Inconsistent background.}\enspace
For \(\mathcal P_{\rm inconsistent}\), no valuation satisfies both
\(a\) and \(\neg a\). Thus neither candidate, \(\varnothing\)
or \(\{a\}\), is consistent with \(T\). There is no explanation
and both selected subfamilies are empty. For each candidate the entailment
of \(m\) alone is true vacuously, since its premise set has no model.

\paragraph{Observation already given.}
The only candidates for \(\mathcal P_{\rm given}\) are
\(\varnothing\) and \(\{a\}\). Both are consistent with
\(T=\{m\}\): assign \(m\) true and, when needed, \(a\) true.
Every model of either augmented theory makes \(m\) true because it is
already in \(T\). Thus both candidates explain. The empty one is a
proper subset of \(\{a\}\) and has size zero, so it is the only
member of both minimality families.

\paragraph{Empty data.}
For \(\mathcal P_{\rm empty}\), the only candidate is the empty set,
and the only valuation on \(V\) is the empty function. It satisfies
the empty theory and every member of the empty observation set. Thus the
empty candidate is an explanation. Having no competitors, it belongs to
both minimality families. Each of the three families therefore has exactly
one member, the empty set.

\paragraph{Unequal sizes.}
Consider any \(E\subseteq\{a,b,c\}\). Assigning \(a,b,c,m\) all
true satisfies \(T\cup E\), so every candidate is consistent. If
\(a\in E\), every model of \(T\cup E\) must make \(m\) true
by \(a\to m\). If \(b,c\in E\), the same conclusion follows
from \((b\land c)\to m\).

For the converse, suppose \(a\notin E\) and
\(\{b,c\}\nsubseteq E\). Assign an atom among \(a,b,c\) true
exactly when it belongs to \(E\), and assign \(m\) false. Then
\(a\) is false and at least one of \(b,c\) is false. Both antecedents
in \(T\) are false, so both implications hold. This valuation satisfies
\(T\cup E\) but falsifies \(m\). Hence \(E\) is not an
explanation. We have proved the exact criterion
\[
 E\in\operatorname{Expl}(\mathcal P_{\rm sizes})
   \quad\Longleftrightarrow\quad
       a\in E\ \text{ or }\ \{b,c\}\subseteq E.
\]
The four subsets containing \(a\) and the one remaining subset
\(\{b,c\}\) are exactly the five listed explanations. The empty set
does not explain, so \(\{a\}\) is subset-minimal. Every proper
subset of \(\{b,c\}\) omits \(a\) and at least one of \(b,c\),
so \(\{b,c\}\) is also subset-minimal. Each other listed explanation
properly contains \(\{a\}\). Thus there are exactly two subset-minimal
explanations, and neither contains the other. Their sizes are one and two,
so the unique minimum-cardinality explanation is \(\{a\}\).

\paragraph{Consistency constraint.}
The four candidates for \(\mathcal P_{\rm constraint}\) are
\[
 \varnothing,\{a\},\{b\},\{a,b\}.
\]
The candidate
\(\{a\}\) is consistent: choose \(a,m\) true and \(b\) false.
It entails \(m\) with \(T\) by the implication \(a\to m\).
For the empty candidate choose \(a,b,m\) all false; this is a model
of \(T\) which falsifies \(m\). For \(\{b\}\), choose
\(b\) true and \(a,m\) false; this is a model of \(T\cup\{b\}\)
which falsifies \(m\). These candidates are consistent but do not
explain. No valuation satisfies \(a,b\) and \(\neg(a\land b)\),
so \(\{a,b\}\) is inconsistent with \(T\). Therefore
\(\{a\}\) is the sole explanation and the sole member of both
minimality families. Its failed superset verifies the consistency qualification.

\paragraph{Permitted hypotheses.}
The only permitted candidates for \(\mathcal P_{\rm permitted}\) are
\(\varnothing\) and \(\{a\}\). In each case set \(z,m\) false
and set \(a\) true if the candidate requires it. The implication
\(z\to m\) holds, the candidate is satisfied, and \(m\) is false.
Thus both candidates are consistent and non-explanatory, and all three
explanation families are empty.

Separately, setting \(z,m\) true provides a model of \(T\cup\{z\}\),
and every model of that set makes \(m\) true by \(z\to m\).
Nevertheless \(\{z\}\nsubseteq H=\{a\}\). It is therefore not
an explanation for the declared input. Allowing it would require a different
hypothesis set and hence a different problem.
\end{proof}
\end{annexitem}

\begin{annexitem}{annex:p16.abduction-truth-boundary}{thm:p16-abduction-truth-boundary}{A countervaluation to recovering hypotheses from observations}
\begin{proof}
The preceding calculation proves that \(\{a\}\) and \(\{b,c\}\)
are subset-minimal explanations in \(\mathcal P_{\rm sizes}\).
Define a valuation on its four atoms by
\[
 v(a)=v(b)=v(c)=\mathbf{F},\qquad v(m)=\mathbf{T}.
\]
Both antecedents \(a\) and \(b\land c\) are false, so the two
implications in \(T\) are true. The observation \(m\) is also true.
Thus \(v\) satisfies \(T\cup M\), while it falsifies both
\(a\) and \(b\land c\). Extending this valuation to \(\mathsf{Prop}\)
as in \cref{thm:p16-abduction-semantics} supplies countervaluations in
the book's global semantics. This proves both failures of entailment.
\end{proof}
\end{annexitem}

\begin{annexitem}{annex:p15-w02c-p15-027-formula-hypotheses}{thm:p15-w02c-p15-027-formula-hypotheses}{Formulas}
\begin{proof}
For \(S\subseteq\Phi\) put \(E_S=\{x_\varphi:\varphi\in S\}\). We show
that restriction to \(V\) is a bijection
\[
 \rho_S:\operatorname{Mod}_{V'}(T'\cup E_S)\longrightarrow
        \operatorname{Mod}_V(T\cup S).
\]
Let \(v'\) satisfy \(T'\cup E_S\). The formulas of \(T\) and of
\(\Phi\) are over \(V\), so by the coincidence argument in the proof of
\cref{thm:p16-abduction-semantics} they have the same values under
\(v'\) and under its restriction \(v'|_V\). Hence \(v'|_V\) satisfies
\(T\). For \(\varphi\in S\), \(v'\) makes \(x_\varphi\) true and
satisfies \(x_\varphi\leftrightarrow\varphi\), so it makes \(\varphi\)
true, and so does \(v'|_V\). Thus \(\rho_S\) maps into
\(\operatorname{Mod}_V(T\cup S)\). Conversely, given \(v\) in that set,
define \(v'\) on \(V'\) by \(v'|_V=v\) and \(v'(x_\varphi)\) equal to
the value of \(\varphi\) under \(v\), for every \(\varphi\in\Phi\); this
is a function because the new atoms are distinct and lie outside \(V\).
Then \(v'\) satisfies \(T\), each biconditional
\(x_\varphi\leftrightarrow\varphi\) (both sides have the same value), and
each \(x_\varphi\) with \(\varphi\in S\) (as \(v\) satisfies
\(\varphi\)). So \(\rho_S\) is onto. It is one-to-one because every
\(v'\) satisfying \(T'\) gives \(x_\varphi\) the value of \(\varphi\)
under \(v'|_V\), so \(v'\) is determined by its restriction.

\emph{Item 1.} Take \(S=\varnothing\): \(\operatorname{Mod}_{V'}(T')\) is
nonempty exactly when \(\operatorname{Mod}_V(T)\) is.

\emph{Item 2.} By the bijection, \(\operatorname{Mod}_{V'}(T'\cup E_S)\)
is nonempty exactly when \(\operatorname{Mod}_V(T\cup S)\) is. Every
\(m\in M\) lies in \(V\), so its value under \(v'\) equals its value
under \(\rho_S(v')\). Hence every model of \(T'\cup E_S\) satisfies
every \(m\in M\) exactly when every model of \(T\cup S\) does. By the
finite criterion of \cref{thm:p16-abduction-semantics}, applied to
\(\mathcal P'\), and the same criterion over \(V\) for formula
explanations, \(E_S\in\operatorname{Expl}(\mathcal P')\) exactly when
\(S\) is a formula explanation. The map \(S\mapsto E_S\) is a bijection
from subsets of \(\Phi\) onto subsets of \(H'\), since the atoms
\(x_\varphi\) are distinct; it satisfies \(S_1\subseteq S_2\) exactly
when \(E_{S_1}\subseteq E_{S_2}\), and \(|E_S|=|S|\). It therefore
restricts to a bijection between formula explanations and
\(\operatorname{Expl}(\mathcal P')\). Subset-minimality and minimum
cardinality are defined from inclusion and cardinality alone, so the
bijection carries each minimality family of formula explanations onto
the corresponding family of \cref{def:p16-abduction-minimality}.
\end{proof}
\end{annexitem}

\begin{annexitem}{annex:p15-w02c-p15-027-source-extension}{thm:p15-w02c-p15-027-source-extension}{Adding formula hypotheses beside the atoms of \texorpdfstring{\(H\)}{H}}
\begin{proof}
Write \(X=X_\Phi\) and \(V^{\Phi}=V\cup X\). For \(S\subseteq H\) and
\(\Psi\subseteq\Phi\) we first show that restriction to \(V\) is a
bijection
\[
 \rho_{S,\Psi}:\operatorname{Mod}_{V^{\Phi}}(T^{\Phi}\cup S\cup X_\Psi)
   \longrightarrow\operatorname{Mod}_V(T\cup S\cup\Psi).
\]
Let \(w\) satisfy \(T^{\Phi}\cup S\cup X_\Psi\). The formulas of \(T\),
the atoms of \(S\subseteq H\subseteq V\) and the formulas of \(\Phi\) are
all over \(V\), so by the coincidence argument in the proof of
\cref{thm:p16-abduction-semantics} each has the same value under \(w\) and
under \(w|_V\). Hence \(w|_V\) satisfies \(T\) and \(S\). For
\(\varphi\in\Psi\), \(w\) makes \(x'_\varphi\) true and satisfies
\(x'_\varphi\leftrightarrow\varphi\), so it makes \(\varphi\) true, and so
does \(w|_V\). Thus \(\rho_{S,\Psi}\) maps into
\(\operatorname{Mod}_V(T\cup S\cup\Psi)\).

Conversely, let \(v\) satisfy \(T\cup S\cup\Psi\). Define \(w\) on
\(V^{\Phi}\) by \(w|_V=v\) and, for each \(\varphi\in\Phi\),
\(w(x'_\varphi)\) equal to the value of \(\varphi\) under \(v\). This is a
function, because the atoms \(x'_\varphi\) are distinct and lie outside
\(V\). By coincidence, \(w\) satisfies \(T\) and every atom of \(S\); it
satisfies each biconditional \(x'_\varphi\leftrightarrow\varphi\),
because both sides have the value of \(\varphi\) under \(v\); and it makes
\(x'_\varphi\) true for \(\varphi\in\Psi\), because \(v\) satisfies
\(\varphi\). So \(w\) lies in the domain and \(\rho_{S,\Psi}(w)=v\); the
map is onto. It is one-to-one: every \(w\) satisfying \(T^{\Phi}\) gives
each \(x'_\varphi\) the value of \(\varphi\) under \(w|_V\), so \(w\) is
determined by \(w|_V\).

\emph{Item 1.} Take \(S=\varnothing\) and \(\Psi=\varnothing\): the set
\(\operatorname{Mod}_{V^{\Phi}}(T^{\Phi})\) is nonempty exactly when
\(\operatorname{Mod}_V(T)\) is.

\emph{Item 2.} Every \(E\subseteq H\cup X\) has exactly one
representation \(E=S\cup X_\Psi\) with \(S\subseteq H\) and
\(\Psi\subseteq\Phi\), namely \(S=E\cap H\) and
\(\Psi=\{\varphi:x'_\varphi\in E\}\): the sets \(H\subseteq V\) and \(X\)
are disjoint, and distinct formulas have distinct atoms. So
\((S,\Psi)\mapsto S\cup X_\Psi\) is a bijection from all pairs onto
\(2^{H\cup X}\). By the same disjointness,
\(S_1\cup X_{\Psi_1}\subseteq S_2\cup X_{\Psi_2}\) holds exactly when
\(S_1\subseteq S_2\) and \(X_{\Psi_1}\subseteq X_{\Psi_2}\), that is,
\(\Psi_1\subseteq\Psi_2\); and \(|S\cup X_\Psi|=|S|+|X_\Psi|=|S|+|\Psi|\).
It remains to show that \(S\cup X_\Psi\) is an explanation for
\(\mathcal P^{\Phi}\) exactly when \((S,\Psi)\) is a mixed explanation. By
the bijection \(\rho_{S,\Psi}\), the set
\(\operatorname{Mod}_{V^{\Phi}}(T^{\Phi}\cup S\cup X_\Psi)\) is nonempty
exactly when \(\operatorname{Mod}_V(T\cup S\cup\Psi)\) is. Each
\(m\in M\) lies in \(V\), so its value under \(w\) equals its value under
\(\rho_{S,\Psi}(w)\); since \(\rho_{S,\Psi}\) is onto, every model on the
left makes every \(m\in M\) true exactly when every model on the right
does. By the finite criterion of \cref{thm:p16-abduction-semantics},
applied over \(V^{\Phi}\) on the left and over \(V\) on the right, this
is the required equivalence.

\emph{Item 3.} A pair \((S,\varnothing)\) is a mixed explanation exactly
when \(S\in\operatorname{Expl}(\mathcal P)\), since the two conditions are
then the same. With item~2, a set \(E\subseteq H\), which is
\(E\cup X_\varnothing\), lies in \(\operatorname{Expl}(\mathcal P^{\Phi})\)
exactly when it lies in \(\operatorname{Expl}(\mathcal P)\). Every proper
subset of such a set is again a subset of \(H\), so no proper subset of
\(E\) is an explanation for \(\mathcal P^{\Phi}\) exactly when none is an
explanation for \(\mathcal P\). Hence
\(E\in\operatorname{MinExpl}(\mathcal P^{\Phi})\) exactly when
\(E\in\operatorname{MinExpl}(\mathcal P)\).

\emph{Item 4.} For \(\Phi=\{F\}\) the pairs are \((S,\varnothing)\) and
\((S,\{F\})\). The first are mixed explanations exactly when
\(S\in\operatorname{Expl}(\mathcal P)\), by item~3, and the second
exactly when \(S\in\mathcal S_F\), by definition; item~2 gives the
displayed union. Let \(S\in\mathcal S_F\). The proper subsets of
\(S\cup\{x'\}\) are the sets \(S_0\cup\{x'\}\) with \(S_0\subsetneq S\)
and the sets \(S_0\) with \(S_0\subseteq S\). By item~2 a set of the first
kind is an explanation exactly when \(S_0\in\mathcal S_F\), and by item~3
a set of the second kind is one exactly when
\(S_0\in\operatorname{Expl}(\mathcal P)\). So \(S\cup\{x'\}\) is
subset-minimal exactly when no \(S_0\subsetneq S\) lies in
\(\mathcal S_F\) and no \(S_0\subseteq S\) lies in
\(\operatorname{Expl}(\mathcal P)\).

\emph{Item 5.} Here \(V=\{a,b,m\}\), \(H=\{a,b\}\), \(M=\{m\}\) and
\(T=\{(a\land b)\to m\}\). The candidate \(\{a\}\) is consistent with
\(T\) but does not entail \(m\): the valuation making \(a\) true and
\(b\) and \(m\) false satisfies \(T\) and \(a\). The same valuation with
the roles of \(a\) and \(b\) exchanged refutes \(\{b\}\), and the
valuation making all three atoms false refutes \(\varnothing\). The
candidate \(\{a,b\}\) is consistent (make all three atoms true), and every
model of \(T\) in which \(a\) and \(b\) are true makes \(m\) true. So
\(\operatorname{Expl}(\mathcal P)=\{\{a,b\}\}\), and this single
explanation is both subset-minimal and of minimum cardinality. For
\(F=a\land b\), \(\mathcal S_F\) consists of all four subsets of \(H\):
\(T\cup S\cup\{a\land b\}\) is satisfied by the valuation making every atom
true, and each of its models makes \(a\land b\) and hence \(m\) true. By
item~4,
\(\operatorname{Expl}(\mathcal P^{\{F\}})=\{\{a,b\}\}\cup
\{\{x'\},\{a,x'\},\{b,x'\},\{a,b,x'\}\}\). The set \(\{x'\}\) is
subset-minimal: \(\varnothing\) is inclusion-minimal in \(\mathcal S_F\)
and \(\varnothing\notin\operatorname{Expl}(\mathcal P)\). The sets
\(\{a,x'\}\), \(\{b,x'\}\) and \(\{a,b,x'\}\) properly contain the
explanation \(\{x'\}\), and \(\{a,b\}\) is subset-minimal by item~3. So
\(\operatorname{MinExpl}(\mathcal P^{\{F\}})=\{\{a,b\},\{x'\}\}\). The
least size of an explanation is one, attained only by \(\{x'\}\), so
\(\operatorname{CardExpl}(\mathcal P^{\{F\}})=\{\{x'\}\}\).
\end{proof}
\end{annexitem}

\begin{annexitem}{annex:p15-w02c-p15-027-causal-boundary}{thm:p15-w02c-p15-027-causal-boundary}{Two causal models for one abduction input}
\begin{proof}
Both graphs are acyclic: in \(C_1\) every edge ends at \(m\), which has no
outgoing edge; in \(C_2\) the edges go from \(m\) and \(c\) to \(a\) and
\(b\), which have no outgoing edges. By \cref{thm:p16-causal-solution}
each model has exactly one solution at each outcome, computed in
topological order. Write an outcome as \((u_a,u_b,u_c,u_m)\) and a
valuation as \((a,b,c,m)\).

\emph{Item 1.} The models of \(T=\{a\to m,(b\land c)\to m\}\) are the
eight valuations with \(m=\mathbf T\) together with the three with
\(m=\mathbf F\), \(a=\mathbf F\) and \(b\land c\) false: eleven in all.
In \(C_1\) the solution at an outcome is
\((u_a,u_b,u_c,u_a\lor(u_b\land u_c)\lor u_m)\). Its last coordinate is
true whenever \(a\) or \(b\land c\) is, so it satisfies \(T\). Conversely
a model \((a,b,c,m)\) of \(T\) is the solution at the outcome
\((a,b,c,m)\): if \(m\) is true the equation for \(m\) gives true, and if
\(m\) is false then \(a\), \(b\land c\) and \(u_m=m\) are false. In
\(C_2\) the solution is \(m=u_m\), \(c=u_c\), \(a=u_m\land u_a\),
\(b=u_b\land(u_m\lor\neg u_c)\). If \(u_m\) is true it is
\((u_a,u_b,u_c,\mathbf T)\), which ranges over all eight valuations with
\(m\) true. If \(u_m\) is false then \(a\) is false and \(b\) is true only
when \(c\) is false, so \(b\land c\) is false; the three valuations
\((\mathbf F,\mathbf F,\mathbf F,\mathbf F)\),
\((\mathbf F,\mathbf T,\mathbf F,\mathbf F)\) and
\((\mathbf F,\mathbf F,\mathbf T,\mathbf F)\) arise at the outcomes with
\(u_m\) false and \((u_b,u_c)\) equal to \((\mathbf F,\mathbf F)\),
\((\mathbf T,\mathbf F)\) and \((\mathbf F,\mathbf T)\). So in both models
the attained solutions are exactly the eleven models of \(T\), each at
an outcome of mass at least \(1/16\). The three explanation families are
determined by \(V\), \(H\), \(M\) and \(\operatorname{Mod}_V(T)\), by
\cref{thm:p16-abduction-semantics}, so the families of
\cref{thm:p16-abduction-examples} do not depend on which model is
intended; the input contains neither model.

\emph{Item 2.} In \(C_1\) the intervention \(\operatorname{do}(a=\mathbf T)\)
replaces the equation for \(a\) by \(a=\mathbf T\) and keeps the equation
for \(m\) (\cref{def:p16-causal-intervention}), so at every outcome
\(m=\mathbf T\lor(u_b\land u_c)\lor u_m=\mathbf T\). Under
\(\operatorname{do}(b=\mathbf T,c=\mathbf T)\),
\(m=u_a\lor\mathbf T\lor u_m=\mathbf T\).

\emph{Item 3.} In \(C_2\), \(m\) has no parents and is a descendant of
neither \(a\) nor \(\{b,c\}\): \(a\) and \(b\) have no children, and the
only child of \(c\) is \(b\). By item~2 of
\cref{thm:p16-causal-intervention-laws}, under either intervention \(m\)
keeps its value \(u_m\) at every outcome. The outcomes with \(u_m\) false
have total mass \(8/16=1/2\).

\emph{Item 4.} The valuation \((\mathbf F,\mathbf F,\mathbf F,\mathbf T)\)
is a model of \(T\), so by item~1 it is the solution at an outcome of
positive mass in each model. In \(C_1\) it is attained at the outcome
\((\mathbf F,\mathbf F,\mathbf F,\mathbf T)\), where \(m\) is true only
through \(u_m\). In \(C_2\) it is attained at the same outcome, where
\(a=\mathbf T\land\mathbf F\) and \(b=\mathbf F\land(\mathbf T\lor\mathbf T)\)
are false and \(m=u_m\) is true.
\end{proof}
\end{annexitem}

\begin{annexitem}{annex:p15-w02c-p15-027-theory-forms}{thm:p15-w02c-p15-027-theory-forms}{Clausal theories}
\begin{proof}
\emph{Item 1.} For every candidate \(E\subseteq H\),
\[
\begin{aligned}
 \operatorname{Mod}_V(T'\cup E)
   &=\operatorname{Mod}_V(T')\cap\operatorname{Mod}_V(E)\\
   &=\operatorname{Mod}_V(T)\cap\operatorname{Mod}_V(E)
   =\operatorname{Mod}_V(T\cup E),
\end{aligned}
\]
since a valuation satisfies a union of formula sets exactly when it
satisfies each part. By \cref{thm:p16-abduction-semantics}, whether \(E\)
is an explanation depends only on \(M\) and on this set of models, so
\(\operatorname{Expl}(\mathcal P')=\operatorname{Expl}(\mathcal P)\).
The families \(\operatorname{MinExpl}\) and \(\operatorname{CardExpl}\)
are defined from \(\operatorname{Expl}\) by inclusion and cardinality
alone, so they agree as well.

\emph{Item 2.} Each \(C_w\) contains, for every \(p\in V\), exactly one of
the literals \(p,\neg p\), so its literals are pairwise distinct and its
atoms lie in \(V\). We show that \(v\in\mathsf{Val}(V)\) satisfies \(C_w\)
exactly when \(v\ne w\). If \(v=w\), each literal \(\neg p\) of \(C_w\) has
\(w(p)=\mathbf T\) and so is false, and each literal \(p\) of \(C_w\) has
\(w(p)=\mathbf F\); thus no literal is true and the clause is false
(\cref{lem:prop-clause-satisfaction}; when \(V=\varnothing\), \(C_w\) is
the empty clause, which is false). If \(v\ne w\), choose \(p\) with
\(v(p)\ne w(p)\). When \(w(p)=\mathbf T\), the literal \(\neg p\) is in
\(C_w\) and true under \(v\); when \(w(p)=\mathbf F\), the literal \(p\) is
in \(C_w\) and true under \(v\). So \(v\) satisfies \(C_w\).
With \(W\) as in the statement, \(v\in\operatorname{Mod}_V(T^{\rm cl})\) if and only if \(v\ne w\)
for every \(w\in W\), that is, \(v\notin W\), that is,
\(v\in\operatorname{Mod}_V(T)\). The set \(T^{\rm cl}\) is finite because
\(\mathsf{Val}(V)\) is. Item~1 now gives the statement about explanation
families.

\emph{Item 3.} Induction on \(E\). If \(E\) is an atom \(x\in X\), then
\(E_{\phi(X)}\) is \(\topv\) when \(v(x)=\mathbf T\) and \(\bottom\) when
\(v(x)=\mathbf F\); in both cases its value is \(v(x)\). An atom outside
\(X\), and \(\bottom\), are unchanged by the substitution. For a compound
formula the substitution acts on the immediate subformulas, which keep
their values by the induction hypothesis, and the connective applies the
same truth function to them. The abbreviations \(\topv\) and
\(\leftrightarrow\) are expanded before the induction, and replacing an atom
commutes with that expansion, since the expansion copies the atom's
occurrences unchanged.

\emph{Item 4.} Compare truth tables. The implication \(a\to m\) is false
exactly when \(a\) is true and \(m\) false, which is exactly when
\(m\lor\neg a\) is false. Likewise \((b\land c)\to m\) and
\(m\lor\neg b\lor\neg c\) are false exactly when \(b,c\) are true and \(m\)
is false, \(z\to m\) and \(m\lor\neg z\) exactly when \(z\) is true and
\(m\) false, and \(\neg(a\land b)\) and \(\neg a\lor\neg b\) exactly when
\(a,b\) are both true. Each theory is replaced member by member by an
equivalent formula over the same alphabet, so the model sets agree. In
the form \(x_1\lor\cdots\lor x_k\lor\neg x_{k+1}\lor\cdots\lor\neg x_n\),
the clauses \(m\lor\neg a\), \(m\lor\neg b\lor\neg c\) and \(m\lor\neg z\)
have \(k=1\), and \(\neg a\lor\neg b\) has \(k=0\).
\end{proof}
\end{annexitem}

\begin{annexitem}{annex:p15-w02c-p15-027-preference}{thm:p15-w02c-p15-027-preference}{Preferred solutions under the three preferences}
\begin{proof}
\emph{Item 1.} Equality, inclusion and the comparison of cardinalities are
reflexive and transitive, so they are preorders on \(2^H\). For equality,
\(S=S'\) and \(S'\ne S\) cannot both hold, so the strict part is empty.
For inclusion, \(S\subseteq S'\) together with \(S'\nsubseteq S\) says that
\(S\subseteq S'\) and \(S\ne S'\), because \(S\subseteq S'\) and
\(S'\subseteq S\) hold together exactly when \(S=S'\). For the size order,
\(|S|\le|S'|\) together with the failure of \(|S'|\le|S|\) is equivalent to
\(|S|<|S'|\), since the order of the natural numbers is total. Substituting
these strict parts into the definition of \(\operatorname{Sol}_{\preceq}\):
no solution is strictly below anything under equality, so every solution
is kept; a solution is kept under \(\subseteq\) exactly when no solution is
a proper subset of it, which is the definition of
\(\operatorname{MinExpl}\); and it is kept under \(\le\) exactly when no
solution has fewer members, which is \(|E|\le|F|\) for every explanation
\(F\), the definition of \(\operatorname{CardExpl}\).

\emph{Item 2.} \(S\prec S\) would require both \(S\preceq S\) and
\(S\not\preceq S\), so \(\prec\) is irreflexive. If \(S_1\prec S_2\) and
\(S_2\prec S_3\), then \(S_1\preceq S_3\) by transitivity; and if
\(S_3\preceq S_1\) held, then \(S_3\preceq S_1\preceq S_2\) would give
\(S_3\preceq S_2\), contrary to \(S_2\prec S_3\). Hence \(S_1\prec S_3\).
By definition \(\operatorname{Sol}_{\preceq}(\mathcal P)\subseteq
\operatorname{Expl}(\mathcal P)\), so the former is empty when the latter
is. Suppose \(\operatorname{Expl}(\mathcal P)\) is nonempty with \(k\)
members and has no \(\prec\)-minimal member. Starting from any member
\(S_0\), choose successively \(S_{j+1}\in\operatorname{Expl}(\mathcal P)\)
with \(S_{j+1}\prec S_j\), for \(j=0,\ldots,k-1\). Among the \(k+1\) sets
\(S_0,\ldots,S_k\) two coincide, say \(S_i=S_l\) with \(i<l\). By
transitivity along \(S_l\prec S_{l-1}\prec\cdots\prec S_i\) we get
\(S_l\prec S_i=S_l\), contradicting irreflexivity. So a \(\prec\)-minimal
solution exists.

\emph{Item 3.} Let \(S,S'\) be solutions with \(S\subsetneq S'\). Then
\(S\subseteq S'\) and \(S\ne S'\), so \(S\) is strictly below \(S'\) under
\(\subseteq\) by item~1. Since \(S'\) is finite and contains \(S\) and at
least one further element, \(|S'|\ge|S|+1>|S|\), so \(S\) is strictly
below \(S'\) under \(\le\). Under equality nothing is strictly below
anything, so the irredundancy condition holds exactly when there is no
pair of solutions with \(S\subsetneq S'\). For \(\mathcal P_{\rm given}\),
\cref{thm:p16-abduction-examples} gives the solutions \(\varnothing\) and
\(\{a\}\), and \(\varnothing\subsetneq\{a\}\).

\emph{Item 4.} Let \(S\in\operatorname{Sol}_{\preceq}(\mathcal P)\) and
suppose \(S\notin\operatorname{Sol}_{\subseteq}(\mathcal P)\). By item~1
there is a solution \(S'\subsetneq S\). Irredundancy gives \(S'\prec S\),
contradicting the \(\prec\)-minimality of \(S\).
\end{proof}
\end{annexitem}

\begin{annexitem}{annex:p15-w02c-p15-027-decomposition}{thm:p15-w02c-p15-027-decomposition}{Decomposition over disjoint alphabets}
\begin{proof}
\textbf{Step 1: splitting candidates.}\enspace
Since \(H_i\subseteq V_i\) and the \(V_i\) are pairwise disjoint, the
\(H_i\) are pairwise disjoint. For \(S\subseteq H\) put \(S_i=S\cap H_i\).
Then \(S=S_1\cup\cdots\cup S_n\) with pairwise disjoint parts, so
\(|S|=\sum_i|S_i|\). Conversely, if \(S_i\subseteq H_i\) for each \(i\),
the union \(S\) satisfies \(S\cap H_i=S_i\). Hence
\(S\mapsto(S_1,\ldots,S_n)\) is a bijection from \(2^H\) onto the tuples of
subsets of the \(H_i\), and \(S\subseteq S'\) holds exactly when
\(S_i\subseteq S'_i\) for every \(i\). This gives the uniqueness of the
representation.

\textbf{Step 2: gluing valuations.}\enspace
Because the \(V_i\) are pairwise disjoint, valuations
\(v_i\in\mathsf{Val}(V_i)\) combine into one valuation
\(v\in\mathsf{Val}(V)\) with \(v|_{V_i}=v_i\), and every
\(v\in\mathsf{Val}(V)\) is the combination of its restrictions. Each
formula of \(T_i\cup S_i\) is over \(V_i\), so by the coincidence argument
in the proof of \cref{thm:p16-abduction-semantics} it has the same value
under \(v\) and under \(v|_{V_i}\). As
\(T\cup S=\bigcup_i(T_i\cup S_i)\), a valuation \(v\in\mathsf{Val}(V)\)
satisfies
\[
 v\in\operatorname{Mod}_V(T\cup S)
  \quad\text{exactly when}\quad
  v|_{V_i}\in\operatorname{Mod}_{V_i}(T_i\cup S_i)
      \text{ for every }i.
\]
Call this the gluing criterion.

\textbf{Step 3: the case of equality.}\enspace
We show \(S\in\operatorname{Sol}(\mathcal P)\) if and only if
\(S_i\in\operatorname{Sol}(\mathcal P_i)\) for every \(i\).
Suppose every \(S_i\) is a solution. Choose
\(v_i\in\operatorname{Mod}_{V_i}(T_i\cup S_i)\) and combine them; by the
gluing criterion the result is a model of \(T\cup S\), so \(T\cup S\) is
consistent. Let \(m\in M\), say \(m\in M_i\), and let \(v\) be any model of
\(T\cup S\). By the gluing criterion, \(v|_{V_i}\) is a model of
\(T_i\cup S_i\), so it makes \(m\) true, and \(v(m)=v|_{V_i}(m)\) since
\(m\in V_i\). Hence \(T\cup S\entails M\).
Conversely, let \(S\) be a solution and \(v\) a model of \(T\cup S\). By the
gluing criterion each \(v|_{V_i}\) is a model of \(T_i\cup S_i\), so these
sets are consistent. Let \(m\in M_i\) and
\(w_i\in\operatorname{Mod}_{V_i}(T_i\cup S_i)\). Combine \(w_i\) with the
restrictions \(v|_{V_j}\) for \(j\ne i\). By the gluing criterion this is a
model of \(T\cup S\), so it makes \(m\) true; its value at \(m\) is
\(w_i(m)\). Hence \(T_i\cup S_i\entails M_i\).
With step~1 this proves the lemma for equality. If some
\(\operatorname{Sol}(\mathcal P_i)\) is empty, in particular if \(T_i\) is
inconsistent, both sides are empty.

\textbf{Step 4: the subset order.}\enspace
Let \(S\) be a solution, so each \(S_i\) is a solution of
\(\mathcal P_i\). If some \(S_i\) is not subset-minimal, choose a solution
\(S'_i\subsetneq S_i\) of \(\mathcal P_i\) and replace \(S_i\) by
\(S'_i\). By step~3 the resulting set is a solution of \(\mathcal P\), and
by step~1 it is a proper subset of \(S\); so \(S\) is not subset-minimal.
If instead \(S\) is not subset-minimal, take a solution \(S'\subsetneq S\).
By steps~1 and~3, every \(S'_j\) is a solution of \(\mathcal P_j\) with
\(S'_j\subseteq S_j\), and \(S'_j\ne S_j\) for some \(j\), since otherwise
\(S'=S\); so that \(S_j\) is not subset-minimal. Therefore \(S\) is
subset-minimal exactly when every \(S_i\) is, which with step~3 and
\cref{thm:p15-w02c-p15-027-preference} proves the lemma for \(\subseteq\).

\textbf{Step 5: the size order.}\enspace
Let \(S\) be a solution. If every \(S_i\) has minimum size among the
solutions of \(\mathcal P_i\), then every solution \(S'\) of \(\mathcal P\)
has solutions \(S'_i\) as parts, by step~3, and
\(|S'|=\sum_i|S'_i|\ge\sum_i|S_i|=|S|\). If some \(S_i\) does not have
minimum size, replace it by a solution \(S'_i\) of \(\mathcal P_i\) with
\(|S'_i|<|S_i|\); the result is a solution of \(\mathcal P\) with fewer
members than \(S\). Therefore \(S\) has minimum size exactly when every
\(S_i\) has, which with step~3 and
\cref{thm:p15-w02c-p15-027-preference} proves the lemma for \(\le\).

\textbf{Step 6: the overlapping example.}\enspace
For \(\mathcal P_1\), the empty candidate fails, since the valuation with
\(a,m\) false satisfies \(a\to m\) and falsifies \(m\); the candidate
\(\{a\}\) is consistent (make \(a,m\) true) and every model of
\(\{a\to m,a\}\) makes \(m\) true. So
\(\operatorname{Sol}(\mathcal P_1)=\{\{a\}\}\). For \(\mathcal P_2\), the
only candidate is \(\varnothing\), \(\{\neg a\}\) is consistent and there
are no observations, so \(\operatorname{Sol}(\mathcal P_2)=\{\varnothing\}\).
Each family has one member, so it is also its own subset-minimal and
minimum-size family, and the right side is \(\{\{a\}\}\) for all three
preferences. The union input is
\((\{a,m\},\{a\},\{m\},\{a\to m,\neg a\})\). Its empty candidate fails by
the valuation with \(a,m\) false, and \(\{a\}\) is inconsistent with
\(\neg a\). So the union has no solution, and the left side is empty
for all three preferences.
\end{proof}
\end{annexitem}

\begin{annexitem}{annex:p15-w02c-p15-027-relevance}{thm:p15-w02c-p15-027-relevance}{Relevance, necessity and dispensability}
\begin{proof}
\emph{Item 1.} Dispensability is defined as the negation of necessity.

\emph{Item 2.} This is item~2 of \cref{thm:p15-w02c-p15-027-preference}.

\emph{Item 3.} Every \(h\in H\) is either not \(\preceq\)-relevant, in
which case \(h\in I\), or \(\preceq\)-relevant, in which case \(h\in N\) or
\(h\in R\) according as it is \(\preceq\)-necessary or not. So
\(I\cup R\cup N=H\). The sets \(I\) and \(R\) are disjoint because members
of \(R\) are relevant, and \(R\) and \(N\) are disjoint because members of
\(R\) are not necessary. Let \(h\in N\) and choose
\(S\in\operatorname{Sol}_{\preceq}(\mathcal P)\), which exists by
hypothesis. Necessity gives \(h\in S\), so \(h\) is relevant and
\(h\notin I\). Hence \(I\cap N=\varnothing\).

\emph{Item 4.} With no acceptable solution, no \(h\) lies in one, so every
\(h\) is \(\preceq\)-irrelevant; and the requirement that \(h\) lie in every
acceptable solution holds vacuously, so every \(h\) is
\(\preceq\)-necessary. Thus \(I=N=H\), and \(R=\varnothing\) since its
members would have to be relevant. If \(H\ne\varnothing\), then
\(I\cap N=H\ne\varnothing\). For \(\mathcal P_{\rm permitted}\),
\cref{thm:p16-abduction-examples} gives
\(\operatorname{Expl}(\mathcal P_{\rm permitted})=\varnothing\), so by
item~2 there is no acceptable solution under any preorder, and
\(H=\{a\}\ne\varnothing\).

\emph{Item 5.} By \cref{thm:p16-abduction-examples} and item~1 of
\cref{thm:p15-w02c-p15-027-preference},
\(\operatorname{Sol}_{=}(\mathcal P_{\rm sizes})\) consists of
\(\{a\},\{a,b\},\{a,c\},\{a,b,c\},\{b,c\}\),
\(\operatorname{Sol}_{\subseteq}(\mathcal P_{\rm sizes})=\{\{a\},\{b,c\}\}\),
and \(\operatorname{Sol}_{\le}(\mathcal P_{\rm sizes})=\{\{a\}\}\). In the
first two families, \(a\) lies in \(\{a\}\) and \(b,c\) lie in \(\{b,c\}\),
so all three are relevant; \(a\notin\{b,c\}\) and \(b,c\notin\{a\}\), so
none is necessary. In the third family the only acceptable solution is
\(\{a\}\), which contains \(a\) and neither \(b\) nor \(c\).
\end{proof}
\end{annexitem}
